%%%%%%%%%%%%%%%%%%%%%%%%%%%%%%%%%%%%%%%%%
% Version 1.5.2 (2019-06-12)
%
% Created by Vel (vel@latextemplates.com) and Johannes Böttcher
% Based on https://www.latextemplates.com/template/masters-doctoral-thesis
%
% Modified by Nils Siegle & Christian Hauff
%	
% Template license:
% CC BY-NC-SA 4.0 (https://creativecommons.org/licenses/by-nc-sa/4.0/)
%%%%%%%%%%%%%%%%%%%%%%%%%%%%%%%%%%%%%%%%%


%----------------------------------------------------------------------------------------
%	BASIC DOCUMENT CONFIGURATIONS
%----------------------------------------------------------------------------------------

\documentclass[
11pt, % The default document font size, options: 10pt, 11pt, 12pt
oneside, % One side output by default, comment it out to switch to two side with alternating margins
ngerman, % ngerman for (new) German -> (New spelling, important for line breaks in a word)
singlespacing, % Single line spacing, alternatives: onehalfspacing or doublespacing
% draft, % Uncomment to enable draft mode (no pictures, no links, overfull hboxes indicated)
% nolistspacing, % If the document is onehalfspacing or doublespacing, uncomment this to set spacing in lists to single
liststotoc, % Uncomment to add the list of figures/tables/etc to the table of contents
% toctotoc, % Uncomment to add the main table of contents to the table of contents
% parskip, % Uncomment to add space between paragraphs
% nohyperref, % Uncomment to not load the hyperref package
headsepline, % Uncomment to get a line under the header
chapterinoneline, % Uncomment to place the chapter title next to the number on one line
% consistentlayout, % Uncomment to change the layout of the declaration, abstract and acknowledgements pages to match the default layout
]{MastersDoctoralThesis} % The class file specifying the document structure


%----------------------------------------------------------------------------------------
%	MARGIN SETTINGS
%----------------------------------------------------------------------------------------

\geometry{
	paper=a4paper, % Change to letterpaper for US letter
	inner=2.8cm, % Inner margin
	outer=2.8cm, % Outer margin
	bindingoffset=.5cm, % Binding offset
	top=1.5cm, % Top margin
	bottom=1.5cm, % Bottom margin
	% showframe, % Uncomment to show how the type block is set on the page
}


%----------------------------------------------------------------------------------------
%	PACKAGES
%----------------------------------------------------------------------------------------

% ------ FONTS ------
\usepackage[T1]{fontenc} % Output font encoding
\usepackage[utf8]{inputenc} % Required for inputting international characters
\usepackage{lmodern,textcomp} %Latin modern font
\usepackage{mathpazo} % Use the Palatino font by default

% ------ BIBLIOGRAPHY ------
\usepackage[backend=bibtex,style=authoryear,natbib=true]{biblatex} % Use the bibtex backend with the authoryear citation style (which resembles APA); DIN ISO: science
\addbibresource{sources.bib} % The filename of the bibliography

% ------ FIGURES ------
\usepackage{graphicx}
\usepackage[font=small,labelfont=bf]{caption} % Required for specifying captions to tables and figures
\usepackage{subcaption} % Required for nested captions
\usepackage{wrapfig} % Required for wrapping text around figures
% \usepackage{forest} % For generating directory trees

% ------ TRANSLATIONS ------
\usepackage{translator}

% Include translations
% ---------- CUSTOM TRANSLATIONS ----------

% Define translations for "Chapter"
\newtranslation[to=German]{chapter}
  {Kapitel}
\newtranslation[to=English]{chapter}
  {chapter}

% Define translations for "Figure"
\newtranslation[to=German]{figure}
  {Abbildung}
\newtranslation[to=English]{figure}
  {figure}

% Define translations for "Appendix"
\newtranslation[to=German]{appendix}
  {Anhang}
\newtranslation[to=English]{appendix}
  {appendix}

% ---------- TITLE PAGE ----------

% Define translations for "Study Programme"
\newtranslation[to=German]{studyProgramme}
  {Studiengang}
\newtranslation[to=English]{studyProgramme}
  {Study Programme}

% Define translations for "Submitted by"
\newtranslation[to=German]{submittedBy}
  {Vorgelegt von}
\newtranslation[to=English]{submittedBy}
  {Submitted by}

% Define translations for "Matriculation Number"
\newtranslation[to=German]{matNum}
  {Matrikelnummer}
\newtranslation[to=English]{matNum}
  {Matriculation Number}

% Define translations for "Initial Examiner"
\newtranslation[to=German]{initEx}
  {Erstprüfer}
\newtranslation[to=English]{initEx}
  {Initial Examiner}

% Define translations for "Secondary Examiner"
\newtranslation[to=German]{secEx}
  {Zweitprüfer}
\newtranslation[to=English]{secEx}
  {Secondary Examiner}

% ---------- CUSTOM TRANSLATIONS ----------

% Use \translate{idToPrint} to print the correct translation
% Define translations for ""
% \newtranslation[to=German]{idToPrint}
%   {German Translation}
% \newtranslation[to=English]{idToPrint}
%   {English Translation}

% ------ OTHER ------
\usepackage{shortvrb}
\PassOptionsToPackage{hyphens}{url}\usepackage{hyperref}

\usepackage{float}
\restylefloat{table}

\usepackage[autostyle=true]{csquotes} % Required to generate language-dependent quotes in the bibliography; generates quotes with \enquote{Text}
\hbadness=99999 %suppresses most underfull hbox warnings

% ------ LISTINGS ------
\usepackage{listings}

% Define JavaScript language support
\lstdefinelanguage{JavaScript}{
  keywords={break, case, catch, continue, debugger, default, delete, do, else, false, finally, for, function, if, in, instanceof, new, null, return, switch, this, throw, true, try, typeof, var, void, while, with},
  morecomment=[l]{//},
  morecomment=[s]{/*}{*/},
  morestring=[b]',
  morestring=[b]",
  sensitive=true
}

% Load required languages
\lstloadlanguages{PHP, XML, HTML, Java, JavaScript, SQL}

% Define colors for highlighting
\definecolor{lstbg}{gray}{0.95}
\definecolor{lstComment}{RGB}{51, 102, 0}
\definecolor{lstKey}{RGB}{0, 51, 204}
\definecolor{lstStr}{RGB}{162, 43, 43}

% Styling
\lstset{
	showtabs=false,
	showspaces=false,
	%tabsize=4,
	%tab=\rightarrowfill,
	showstringspaces=false,
	basicstyle=\ttfamily,
	columns=fullflexible,
	frame=single,
	%frameround=tttt,
	breaklines=true,
	backgroundcolor=\color{lstbg},
	commentstyle=\color{lstComment},
	stringstyle=\color{lstStr},
	keywordstyle=\color{lstKey},
	%prebreak=\mbox{\textbackslash}%,
}


%----------------------------------------------------------------------------------------
%	ADDITIONAL COMMANDS
%----------------------------------------------------------------------------------------

% Option to move href url to footnote for printing
\newcommand{\printurl}[2]{ % Options below are mutually exclusive!
   \href{#1}{#2} % Uncomment for standard href
  % #2\footnote{#1} % Uncomment for print mode, URL in footnote
}

% Better highlighting for inline code
\newcommand{\linecode}[1]{
  \textbf{\texttt{#1}}
}

% Outputs 'Chapter X' (or translation) as a clickable link
\newcommand{\chapref}[1]{
  \hyperref[#1]{\translate{chapter} \ref{#1}}
}

% Outputs 'Appendix X' (or translation) as a clickable link
\newcommand{\appref}[1]{
  \hyperref[#1]{\translate{appendix} \ref{#1}}
}

%----------------------------------------------------------------------------------------
%	THESIS INFORMATION
%----------------------------------------------------------------------------------------

\thesistitle{Catchy Title} % Your thesis title, print it elsewhere with \ttitle
\supervisor{Prof. Supervisor} % Your supervisor's name, print it elsewhere with \supname
\examiner{Prof. Dr. Examiner} % Your examiner's name, print it elsewhere with \examname
\degree{Bachelor Thesis} % Your degree name, print it elsewhere with \degreename
\author{Max \textsc{Mustermann}} % Your name, print it elsewhere with \authorname
\addresses{} % Your address, print it elsewhere with \addressname

\subject{Computer Sciences} % Your subject area, this is not currently used anywhere in the template, print it elsewhere with \subjectname
\university{\href{https://myserver.com/}{University}} % Your university's name and URL, this is used in the title page and abstract, print it elsewhere with \univname
\faculty{\href{https://myserver.com/}{Faculty}} % Your faculty's name and URL, this is used in the title page and abstract, print it elsewhere with \facname


% ------ HYPERSETUP ------
\AtBeginDocument{
\hypersetup{pdftitle=\ttitle} % Set the PDF's title to your title
\hypersetup{pdfauthor=\authorname} % Set the PDF's author to your name
\hypersetup{hidelinks=true,} % Prints all links black; comment out for default LaTeX behavior
}

\begin{document}

% ------ PAGE NUMBERING ------
\frontmatter % Use roman page numbering style (I, II, III, IV...) for the pre-content pages
\pagenumbering{Roman} % Capital roman numbers, use "roman" or comment for lower case (i, ii, iii, iv...)

\pagestyle{plain} % Default to the plain heading style until the thesis style is called for the body content

\sloppy


%----------------------------------------------------------------------------------------
%	TITLE PAGE
%----------------------------------------------------------------------------------------

\begin{titlepage}
 \begin{center}

  \vspace*{.04\textheight}
  {\scshape\LARGE \univname\par}\vspace{0.8cm} % University name
  {\scshape\LARGE \facname\par}\vspace{0.8cm} % Faculty name
  \textsc{\Large \translate{studyProgramme}:}\\[0.5cm] % Thesis type
  \textsc{\Large \degreename}\\[0.5cm] % Thesis type

  \HRule \\[0.4cm] % Horizontal line
  {\huge \bfseries \ttitle\par}\vspace{0.4cm} % Thesis title
  \HRule \\[1.5cm] % Horizontal line

  \begin{minipage}[t]{0.4\textwidth}
    \begin{flushleft} \large
     \emph{\translate{submittedBy}:}
     \newline
     \authorname
     \newline
     \translate{matNum}: 123456
     \newline
     \addressname
    \end{flushleft}
   \end{minipage}
   \begin{minipage}[t]{0.4\textwidth}
    \begin{flushright} \large
     \emph{\translate{initEx}:} \\
     \href{https://myserver.com}{\supname}\\
     \emph{\translate{secEx}:} \\
     \href{https://myserver.com}{\examname}
    \end{flushright}
   \end{minipage}\\[1cm]

  \vfill

  \vfill

  {\large \today}\\[3cm] % Date
  \includegraphics[width=0.60\columnwidth]{Figures/University_Logo}

  \vfill
 \end{center}
\end{titlepage}


%----------------------------------------------------------------------------------------
%	DECLARATION PAGE (Declaration of Independence)
%----------------------------------------------------------------------------------------

\begin{declaration}
 \addchaptertocentry{\authorshipname} % Add the declaration to the table of contents
 \noindent Ich, \authorname, erkläre hiermit, die vorliegende Arbeit \enquote{\ttitle} selbst angefertigt zu haben.

 \begin{itemize}
  \item Wo ich die veröffentlichten Arbeiten anderer zu Rate gezogen haben, wird dies immer angegeben.
  \item Wo ich aus der Arbeit anderer zitiert haben, ist die Quelle immer angegeben. Abgesehen von solchen Zitaten ist diese Arbeit ganz und gar meine eigene.
  \item Ich habe alle verwendeten Hilfsquellen aufgeführt.\\
 \end{itemize}


 \noindent \authorname:\\
 \rule[0.5em]{25em}{0.5pt} % This prints a line for the signature

 \noindent Datum:\\
 \rule[0.5em]{25em}{0.5pt} % This prints a line to write the date
 \\
 \\

\end{declaration}

\cleardoublepage


%----------------------------------------------------------------------------------------
%	QUOTATION PAGE (For a deep quote about the topic)
%----------------------------------------------------------------------------------------

\vspace*{0.2\textheight}

\noindent\enquote{\itshape Programming today is a race between software engineers striving to build bigger and better idiot-proof programs, and the universe trying to build bigger and better idiots. So far, the universe is winning.}\bigbreak

\hfill Rick Cook


%----------------------------------------------------------------------------------------
%	ABSTRACT PAGE
%----------------------------------------------------------------------------------------
\MakeShortVerb{\|}

\begin{abstract}
 \addchaptertocentry{\abstractname} % Add the abstract to the table of contents
 Hier kommt eine ganz tolle Zusammenfassung / Abstract hin.\\
 \\
 \noindent Beispiel für \texttt{inlineCode} mit texttt.\\
 \noindent Beispiel für \linecode{inlineCode} mit linecode.
\end{abstract}


%----------------------------------------------------------------------------------------
%	ACKNOWLEDGEMENTS
%----------------------------------------------------------------------------------------

\begin{acknowledgements}
 \addchaptertocentry{\acknowledgementname} % Add the acknowledgements to the table of contents
 Ich bedanke mich bei \supname \,für die Betreuung dieser Arbeit.\\
\end{acknowledgements}


%----------------------------------------------------------------------------------------
%	LIST OF CONTENTS/FIGURES/TABLES PAGES
%----------------------------------------------------------------------------------------
{\hypersetup{linkcolor=black} % Black links in all lists, curly brackets for local adjustment
  \tableofcontents % Prints the main table of contents

  \listoffigures % Prints the list of figures

  \listoftables % Prints the list of tables

  \lstlistoflistings %Prints the list of (code) listings
  \addchaptertocentry{Listings}
}


%----------------------------------------------------------------------------------------
%	ABBREVIATIONS
%----------------------------------------------------------------------------------------

\begin{abbreviations}{ll} % Include a list of abbreviations (a table of two columns)

 \textbf{Abkürzung} & \textbf{Bedeutung} \\
 \hline
 \textbf{ACL} & \textbf{A}ccess \textbf{C}ontrol \textbf{L}ist\\
 \textbf{AJAX} & \textbf{A}synchronous \textbf{J}avaScript \textbf{A}nd \textbf{X}ML\\
 \textbf{API} & \textbf{A}pplication \textbf{P}rogramming \textbf{I}nterface\\
 \textbf{CSS} & \textbf{C}ascading \textbf{S}tyle \textbf{S}sheets\\
 \textbf{GUI} & \textbf{G}raphical \textbf{U}ser \textbf{I}nterface\\
 \textbf{HTML} & \textbf{H}yper\textbf{t}ext \textbf{M}arkup \textbf{L}anguage\\
 \textbf{JSON} & \textbf{J}ava\textbf{S}cript \textbf{O}bject \textbf{N}otation\\
 \textbf{PHP} & \textbf{PHP} \textbf{H}ypertext \textbf{P}reprocessor\\
 \textbf{SQL} & \textbf{S}tructured \textbf{Q}uery \textbf{L}anguage\\ 
 \textbf{UID} & \textbf{U}ser \textbf{Id}entifier\\
 \textbf{URL} & \textbf{U}niform \textbf{R}esource \textbf{L}ocator\\
 \textbf{UX} & \textbf{U}ser E\textbf{x}perience\\

 % \textbf{} & \textbf{} \textbf{} \textbf{}\\

\end{abbreviations}


%----------------------------------------------------------------------------------------
%	THESIS CONTENT - CHAPTERS
%----------------------------------------------------------------------------------------

\mainmatter % Begin numeric (1,2,3...) page numbering

\pagestyle{thesis} % Return the page headers back to the "thesis" style

% Include the chapters of the thesis as separate files from the Chapters folder
% Uncomment the lines as you write the chapters

% Kapitel 1

\chapter{Einleitung}

\label{Chapter1}

\begin{itemize}
  \item Hierin steht die Begründung der Arbeit, insbesondere der Zweck und die zu lösende(n) Aufgabe(n).
  \item Das Kapitel beginnt mit der Beschreibung der Problemlage.
  \item Hiervon wird die zu lösende Aufgabe abgeleitet.
  \item Ausgangslage und Zielsituation sind so zu konkretisieren und mit Fakten zu belegen (zu quantifizieren), dass am Ende entschieden werden kann, ob und wie die gestellten Fragen beantwortet sind.
  \item Kurzübersicht der Inhalte
\end{itemize}

Erstes Vorkommen von \ac{ajax}, sowie kurz darauf nochmal \ac{ajax}
% Kapitel 2

\chapter{Grundlagen} % Main chapter title

\label{Chapter2}

\begin{itemize}
  \item Hier gehören alle für das weitere Geschehen erforderlichen Grundlagen hin, aber auch nicht mehr, denn:
  \item Die Thesis ist kein Lehrbuch. Hier muss kein Überblick über weit verbreitete Inhalte und Methoden Ihres Themas gegeben werden.
  \item Beschreiben Sie nur das, was Sie im weiteren Verlauf Ihrer Thesis anwenden.
  \item Der fehlende Bezug zwischen Grundlagen und deren Anwendung in der Thesis ist der häufigste Grund für eine suboptimale Bewertung.
  \item Lassen Sie bitte, wenn möglich, das Wort theoretisch weg, denn es kann auch praktische Grundlagen geben.
\end{itemize}

%----------------------------------------------------------------------------------------
%	Abschnitt 1 - Titel
%----------------------------------------------------------------------------------------
\section{Abschnitt 1}

%-----------------------------------
%	Unterpunkt 1.1 - Untertitel
%-----------------------------------
\subsection{Unterabschnitt 1.1}



%-----------------------------------
%	Unterpunkt 1.2 - Untertitel
%-----------------------------------
\subsection{Unterabschnitt 1.2}



%----------------------------------------------------------------------------------------
%	Abschnitt 2 - Titel
%----------------------------------------------------------------------------------------
\section{Abschnitt 2}


%-----------------------------------
%	Unterpunkt 2.1 - Untertitel
%-----------------------------------
\subsection{Unterabschnitt 2.1}


%-----------------------------------
%	Unterpunkt 2.2 - Untertitel
%-----------------------------------
\subsection{Unterabschnitt 2.2}


%----------------------------------------------------------------------------------------
%	Abschnitt 3 - Titel
%----------------------------------------------------------------------------------------
\section{Abschnitt 3}

% Kapitel 3

\chapter{Vorgehen}

\label{Chapter3}

\begin{itemize}
  \item In diesem Kapitel beschreiben Sie, was Sie getan haben, um die Ziele und die Ergebnisse der Arbeit zu erreichen. Hierzu gehört auch wie, womit, warum (zu welchem Zweck) von wem sonst noch und in welcher Reihenfolge die beschriebenen Aktivitäten durchgeführt wurden.
\end{itemize}

%----------------------------------------------------------------------------------------
%	Abschnitt 1 - Titel
%----------------------------------------------------------------------------------------
\section{Abschnitt 1}

%-----------------------------------
%	Unterpunkt 1.1 - Untertitel
%-----------------------------------
\subsection{Unterabschnitt 1.1}



%-----------------------------------
%	Unterpunkt 1.2 - Untertitel
%-----------------------------------
\subsection{Unterabschnitt 1.2}



%----------------------------------------------------------------------------------------
%	Abschnitt 2 - Titel
%----------------------------------------------------------------------------------------
\section{Abschnitt 2}


%-----------------------------------
%	Unterpunkt 2.1 - Untertitel
%-----------------------------------
\subsection{Unterabschnitt 2.1}


%-----------------------------------
%	Unterpunkt 2.2 - Untertitel
%-----------------------------------
\subsection{Unterabschnitt 2.2}


%----------------------------------------------------------------------------------------
%	Abschnitt 3 - Titel
%----------------------------------------------------------------------------------------
\section{Abschnitt 3}

% Kapitel 4

\chapter{Ergebnisse} % Main chapter title

\label{Chapter4}

\begin{itemize}
  \item Hier finden sich alle Ergebnisse.
  \item Die Ergebnisse sind durch Messungen, Formeln, Diagramme, Tabellen und eine logische, für den Leser nachvollziehbare Argumentation zu belegen.
  \item Kein hierfür erforderlicher Beleg (etwa eine Tabelle oder ein Diagramm) wird in den Anhang "`verbannt"'.
  \item Die Kernergebnisse müssen in Abbildungen, Tabellen und Formeln stecken.
  \item Der Leser muss jedes Ergebnis anhand der gegebenen Zahlen, Daten und Fakten (ZDF ) einfach nachvollziehen können, ohne langwierige eigenen Rechnungen.
  \item Im diesem Kapitel dürfen keine grundlegenden Erläuterungen und Begründungen mehr auftauchen.
  \item Wenn alle Ergebnisse beschrieben sind, steht fest, welche Grundlagen benötigt werden und wie das Vorgehen hierfür zu beschreiben ist, damit der Leser es mit Blick auf die Ergebnisse nachvollziehen kann.
\end{itemize}

%----------------------------------------------------------------------------------------
%	Abschnitt 1 - Titel
%----------------------------------------------------------------------------------------
\section{Abschnitt 1}

%-----------------------------------
%	Unterpunkt 1.1 - Untertitel
%-----------------------------------
\subsection{Unterabschnitt 1.1}



%-----------------------------------
%	Unterpunkt 1.2 - Untertitel
%-----------------------------------
\subsection{Unterabschnitt 1.2}



%----------------------------------------------------------------------------------------
%	Abschnitt 2 - Titel
%----------------------------------------------------------------------------------------
\section{Abschnitt 2}


%-----------------------------------
%	Unterpunkt 2.1 - Untertitel
%-----------------------------------
\subsection{Unterabschnitt 2.1}


%-----------------------------------
%	Unterpunkt 2.2 - Untertitel
%-----------------------------------
\subsection{Unterabschnitt 2.2}


%----------------------------------------------------------------------------------------
%	Abschnitt 3 - Titel
%----------------------------------------------------------------------------------------
\section{Abschnitt 3}

% Kapitel X

\chapter{Fazit} % Main chapter title
\label{Chapter5}

\begin{itemize}
  \item Hier wird die gesamte Arbeit auf einer Seite, maximal zwei Seiten zusammengefasst.
  \item Insbesondere werden die zentralen Ergebnisse dargestellt.
  \item Es werden Vorschläge beschrieben zur Lösung der offen gebliebenen Fragen, oder derer, die sich im Laufe der Arbeit ergeben haben.
  \item In diesem Kapitel folgt auch eine kritische Analyse der eigenen Arbeit.
  \item Ein Ausblick auf die Bearbeitung offener Fragen kann gegeben werden.
\end{itemize}



%----------------------------------------------------------------------------------------
%	THESIS CONTENT - APPENDICES
%----------------------------------------------------------------------------------------

\appendix % Cue to tell LaTeX that the following "chapters" are Appendices

% Include the appendices of the thesis as separate files from the Appendices folder
% Uncomment the lines as you write the Appendices

% % Appendix A

\chapter{Formatierung} % Main appendix title

\label{AppendixA} % For referencing this appendix elsewhere, use \ref{AppendixA}

\section{Aufzählungen}

\subsection{Itemize}

\begin{itemize}
  \item Erstes Item
  \item Zweites
  Mit Zeilenumbruch
\end{itemize}

\subsection{Enumerate}

\begin{enumerate}
  \item Punkt mit eins
  \item Und die zwei
\end{enumerate}

\subsection{Description}

\begin{description}
  \item [Punkt 1] Zu Punkt eins lässt sich viel sagen
  \item [1 Label] zum labeln von Dingen
\end{description}

\section{Code}

\begin{lstlisting}[language=XML]
  <role rolename="manager-gui"/>
  <user username="admin" password="MySecretPassword" roles="manager-gui"/>
\end{lstlisting}
% % Appendix A

\chapter{Bilder} % Main appendix title

\label{AppendixB} % For referencing this appendix elsewhere, use \ref{AppendixB}

\begin{figure}[h] % h is the float option, in this case: try to print at this point of the code
    \centering %center the picture
    \includegraphics[width=0.75\columnwidth]{Figures/HS_PF_Logo_Grau}
    \caption[Logo HSPF]{The Logo of Hochschule Pforzheim University\\ Source: \cite{WissArbeiten}} %[short caption (reference in list of figures)]{Long caption}
    \label{fig:DellVC}
\end{figure}

% %AppendixC

\chapter{Quellcode} % Main appendix title

\label{AppendixC}

%\lstinputlisting[caption=index.php,label=index.php,language=PHP]{../Webseite/index.php}



%----------------------------------------------------------------------------------------
%	BIBLIOGRAPHY
%----------------------------------------------------------------------------------------

\printbibliography[heading=bibintoc]

%----------------------------------------------------------------------------------------

\end{document}
